\section{L8 / 8a GPU, CUDA \& OpenCL}

\begin{KR}{Recipe: Size a CUDA / OpenCL Launch}
    one thread per element $\to$ \#threads = \#elements; pick \texttt{threadsPerBlock} (multiple of 32 = warp); grid $=\lceil$elements/block$\rceil$; 2-D \texttt{[R][C]} $\to$ \texttt{threadIdx.x}$\in0..R{-}1$, \texttt{.y}$\in0..C{-}1$.
\end{KR}

\begin{KR}{Recipe: Convert a Loop to a Kernel}
    drop the loop $\to$ \texttt{i=get\_global\_id(0)} (CUDA \texttt{threadIdx/blockIdx}) $\to$ body on element \texttt{i}; work-items must be independent.
\end{KR}

\begin{KR}{Recipe: Map CUDA $\leftrightarrow$ OpenCL}
    thread/block/grid/SM $\equiv$ work-item/work-group/NDRange/CU; host$\to$device$\to$CU$\to$PE.
\end{KR}

\begin{example2}{Ex.\ S4/5 GPU vs CPU}
    How do GPU and CPU differ; thread/loop relationship?
    \tcblower
    CPU = latency-optimised (few fast threads); GPU = \important{throughput} (many slow threads). Each loop iteration $\to$ one GPU thread.
\end{example2}

\begin{example2}{Ex.\ S6/7 Thread \& block organisation}
    Into how many threads is a $[10][5]$ op divided? threadID range?
    \tcblower
    One per element $\to$ \important{$10\times5=50$ threads}; \texttt{threadIdx.x}$\in0..9$, \texttt{threadIdx.y}$\in0..4$ (IDs $(0,0)..(9,4)$). Threads grouped into blocks, blocks into a grid.
\end{example2}

\begin{example2}{Ex.\ S8/9 Execution model / SM}
    What is an SM? How are blocks executed?
    \tcblower
    \important{SM} = Streaming Multiprocessor; a block maps to an SM, which runs threads in \important{warps of 32}. The grid of blocks is distributed across SMs (queue the rest). Size blocks $\times$32.
\end{example2}

\begin{example2}{Ex.\ S10/11 Limitations}
    Why not one thread per pixel for everything?
    \tcblower
    Divergent branches serialise a warp; memory coalescing matters; too-fine work wastes scheduling. Best with regular, independent, arithmetic-heavy data.
\end{example2}

\begin{example2}{Ex.\ S14/15 Rendering pipeline}
    Name the stages; why is it so parallel?
    \tcblower
    world coords $\to$ vertex $\to$ assembly $\to$ (geometry) $\to$ rasterise $\to$ fragment (per pixel) $\to$ blend. Data \important{expands} few vertices $\to$ millions of pixels = massively parallel per-pixel work.
\end{example2}

\begin{example2}{Ex.\ S21--26 OpenCL models}
    Name the four models and the memory regions.
    \tcblower
    Platform / Execution / Memory / Programming. work-item/work-group/NDRange; \important{weak consistency}; regions Global / Constant / Local / Private.
\end{example2}

\begin{example2}{Ex.\ S27/28 Loop $\to$ OpenCL kernel}
    Rewrite \texttt{for(i<4096) c[i]=a[i]+b[i];} and give the NDRange.
    \tcblower
\begin{lstlisting}[language=C, style=basesmol, aboveskip=-0.05\baselineskip, belowskip=-0.5\baselineskip]
__kernel void vadd(__global float*a,*b,*c){
  int i=get_global_id(0); c[i]=a[i]+b[i]; }
\end{lstlisting}
    NDRange = 4096 work-items (e.g.\ work-group 512); loop replaced by the index space.
\end{example2}

\begin{example2}{Ex.\ S29 Coarse vs fine scheduling}
    Two scheduling levels; why interleave warps?
    \tcblower
    \important{Coarse}: blocks $\to$ SMs. \important{Fine}: warp scheduler issues 32-thread warps, dual-issue, \emph{interleaves} when one stalls on memory, another runs $\to$ \important{hides latency}.
\end{example2}

\begin{example2}{Ex.\ S31--33 OpenCL queue \& sync}
    Command-queue ordering? Where can you synchronise?
    \tcblower
    Host enqueues in order; completion in/out-of-order. Sync \important{only within a work-group} (barrier, fence, atomics) none between work-groups.
\end{example2}

\begin{example2}{Ex.\ S40/41 AMP recap}
    Different ISA under one OS?
    \tcblower
    \important{No} separate OSes coupled by openAMP (\texttt{remoteproc} + \texttt{rpmsg}/\texttt{virtio}). Same as L7.
\end{example2}
